% CoRe-TFM manuscript for Machine Learning (Springer Nature)
% Evidence-audited reliability-aware revision.
\documentclass[pdflatex,sn-basic]{sn-jnl}
\usepackage{graphicx}
\usepackage{amsmath,amssymb,amsfonts}
\usepackage{amsthm}
\usepackage{booktabs}
\usepackage{xcolor}
\usepackage{url}
\newcommand{\TV}{\operatorname{TV}}
\newcommand{\KL}{\operatorname{KL}}
\newcommand{\E}{\mathbb{E}}
\newcommand{\Jone}{J^{B\rightarrow A}}
\newcommand{\Jtwo}{J^{A\rightarrow B}}
\newcommand{\hardcore}{Q_{\mathrm{hard}}}
\newcommand{\softcore}{Q_{\mathrm{soft}}}
\theoremstyle{plain}
\newtheorem{proposition}{Proposition}

\begin{document}
\title[CoRe-TFM]{Coherence Is Not Accuracy: Reliability-Aware Reconciliation of Tabular Foundation Model Predictions}

\author*[1]{\fnm{Badampudi Agasthya} \sur{Anirudh}}\email{badampudi.agasthya2023@vitstudent.ac.in}
\author[1]{\fnm{Harshita} \sur{Bogineni}}\email{harshita.bogineni2023@vitstudent.ac.in}
\affil*[1]{\orgdiv{Integrated M.Tech Data Science}, \orgname{[Institution to be completed]}, \orgaddress{\country{India}}}

\abstract{Tabular foundation models (TFMs) can expose direct marginal and conditional probabilities that are not representable by one common joint distribution. We study post-hoc reconciliation of four views without retraining the underlying model. Hard CoRe exactly preserves direct marginals, Soft CoRe relaxes marginal preservation through KL penalties, and Selective CoRe uses held-out labels to choose among raw factorizations, pooling, and reconciliation. Controlled known-truth experiments show that exact marginal preservation helps when direct marginals are reliable and can be strongly harmful when they are not. Across 100 heterogeneous controlled tasks, Selective CoRe improves over arithmetic pooling by $0.006660$ mean exact NLL, winning 68/100 tasks ($p=1.25\times10^{-11}$), with mean full-family oracle regret $0.000710$. The conclusion reverses in a bounded real-model benchmark over ten data sets, five folds, and two released TFMs: Selective CoRe is worse than arithmetic pooling by $0.001831$ joint NLL (95\% dataset-bootstrap CI $[0.000542,0.003291]$; Wilcoxon $p=0.02734$), while reducing marginal distortion. Effects reverse by upstream model. Compatibility, calibration, and predictive correctness are therefore distinct properties, and reconciliation should be treated as a reliability-dependent decision rather than a mandatory projection step.}

\keywords{tabular foundation models, probabilistic consistency, probability reconciliation, reliability, proper scoring rules, Sinkhorn scaling}
\maketitle

\section{Introduction}\label{sec:intro}
Tabular foundation models (TFMs) increasingly provide probability predictions that are consumed beyond ordinary single-target classification. Conditional queries, scenario generation, uncertainty-sensitive decisions, and multivariate simulation can require several probability views from the same frozen model. For two categorical targets $A$ and $B$, consider direct marginals $p(A\mid X)$ and $p(B\mid X)$ and conditionals $p(A\mid B,X)$ and $p(B\mid A,X)$. These define two chain-rule joints,
\begin{align}
\Jone(a,b)&=p_B(b)p_{A\mid B}(a\mid b),\\
\Jtwo(a,b)&=p_A(a)p_{B\mid A}(b\mid a).
\end{align}
Compatibility requires these joints and their implied marginals to agree. Recent work shows that learned conditional systems, including TFMs, need not satisfy such requirements \citep{klotergens2026agree,majid2025consistent}.

The central difficulty is that mathematical coherence is easy to manufacture: any normalized joint is internally coherent by construction. Consequently, zero post-repair inconsistency does not establish better calibration or predictive accuracy. A coherent repair can worsen proper scores or distort useful dependence. CoRe-TFM studies this distinction directly.

We formulate reconciliation as a family of information projections. Arithmetic and geometric pools provide KL-barycenter baselines. Hard CoRe projects a consensus onto the transportation polytope defined by the direct marginals. Soft CoRe replaces exact marginal constraints with KL penalties. Selective CoRe treats reconciliation as optional and chooses among raw factorizations, pools, and repairs using held-out labels.

The paper therefore asks when reconciliation is statistically justified. Controlled known-truth experiments identify the trust assumptions under which repair helps or hurts. A heterogeneous 100-task mixture tests adaptive selection under varying view reliability. A completed bounded-context real benchmark evaluates two released TFMs across ten data sets and shows that the large selector can lose to simple pooling under limited validation.

\section{Related work}\label{sec:related}
TabPFN established prior-data fitted networks as pretrained in-context tabular predictors \citep{hollmann2025tabpfn}; more recent systems such as TabPFN-3 and TabICLv2 extend the paradigm in scale and inference design \citep{grinsztajn2026tabpfn3,qu2026tabiclv2}. Proper-scoring and uncertainty studies emphasize that strong classification accuracy need not imply strong probabilistic reliability \citep{landsgesell2026scoringbench,johnson2026uq}.

The closest diagnostic work formalizes marginalization and factorization consistency for TFM outputs \citep{klotergens2026agree}. Compatibility of conditionals is older still; incompatible conditional families can define compromise distributions without corresponding to one exact joint \citep{mure2017compromise}. CoRe-TFM does not claim to introduce conditional compatibility. It studies post-hoc reconciliation and the associated reliability trade-off for frozen predictors.

Hard CoRe uses classical information projection and iterative proportional fitting \citep{benamou2015bregman}. Generalized scaling for relaxed marginal constraints is related to unbalanced optimal-transport methods \citep{cuturi2013sinkhorn,chizat2018scaling}. These are established optimization tools used here for a TFM-specific reconciliation problem.

\section{Problem formulation}\label{sec:problem}
For categorical targets, factorization inconsistency is
\begin{equation}
\Delta_{\mathrm{fac}}(x)=\TV(\Jone,\Jtwo)=\frac12\sum_{a,b}|\Jone(a,b)-\Jtwo(a,b)|.
\end{equation}
Each factorization automatically preserves one direct marginal and may violate the other. Reconciliation seeks a normalized joint $Q$ that balances factorization agreement, marginal fidelity, and predictive quality.

The four probability views are overlapping queries from one predictor, not independent experts. The objectives should therefore be interpreted as regularized compatibility criteria rather than Bayesian fusion of independent evidence.

\section{Methods}\label{sec:methods}
\subsection{Pooling baselines}
For $w\in[0,1]$,
\begin{align}
Q_{\mathrm{arith},w}&=w\Jone+(1-w)\Jtwo,\\
M_w(a,b)&=\frac{\Jone(a,b)^w\Jtwo(a,b)^{1-w}}{\sum_{a',b'}\Jone(a',b')^w\Jtwo(a',b')^{1-w}}.
\end{align}
Arithmetic and geometric pooling are forward- and reverse-KL barycenters, respectively.

\subsection{Hard CoRe}
Let
\begin{equation}
\mathcal{T}(p_A,p_B)=\{Q\ge0:\sum_bQ(a,b)=p_A(a),\;\sum_aQ(a,b)=p_B(b)\}.
\end{equation}
Hard CoRe solves
\begin{equation}
\hardcore=\arg\min_{Q\in\mathcal{T}(p_A,p_B)}\KL(Q\Vert M_w).
\end{equation}
The feasible set is nonempty and, for positive $M_w$, strict convexity yields a unique optimum.

\subsection{Soft CoRe}
Soft CoRe relaxes exact marginal preservation:
\begin{align}
\softcore=\arg\min_Q\;&w\KL(Q\Vert\Jone)+(1-w)\KL(Q\Vert\Jtwo)\nonumber\\
&+\lambda_A\KL(Q_A\Vert p_A)+\lambda_B\KL(Q_B\Vert p_B).
\end{align}
The KL chain rule decomposes each factorization term into a marginal and conditional contribution, making this a genuine four-view objective. Generalized Sinkhorn-style updates use $\tau_A=\lambda_A/(1+\lambda_A)$ and $\tau_B=\lambda_B/(1+\lambda_B)$.

\subsection{Selective CoRe}
The validation-selected family contains the two raw factorization orders, arithmetic pooling, weighted geometric pooling, weighted Hard CoRe, and weighted Soft CoRe. Five direction weights and seven marginal penalties yield 48 distinct candidates after duplicate removal.

For probability floor $\epsilon$, bounded log loss $B=-\log\epsilon$, validation size $n_v$, and candidate-family size $M$, finite-class concentration gives
\begin{equation}
L(\hat Q)\le\min_{Q\in\mathcal{C}}L(Q)+2B\sqrt{\frac{\log(2M/\delta)}{2n_v}}.
\end{equation}
The qualitative implication is that selection complexity must be matched to validation evidence.

\section{Experimental design}\label{sec:experiments}
We use four evidence layers. First, a surrogate multiclass DGP provides exact $P^*(A,B\mid X)$. Second, an exact-view perturbation benchmark independently perturbs direct marginals and conditional families. Third, a 100-task heterogeneous mixture varies view reliability independently and evaluates Selective CoRe against a full-family oracle. Fourth, a validation-size study evaluates finite-validation selection regret.

The real benchmark contains Anneal, Car, Credit, Customer, Diamonds, Marketing, MIC, Nursery, Phishing, and Wine. Each of five deterministic outer folds uses at most 256 training rows, 52 disjoint validation rows, and 128 test rows. Primary TFMs are TabICLv2 2.1.1 and TabPFN 8.4.0 through the TabPFN-3 interface. CatBoost 1.2.10 is included only as a non-TFM boundary baseline. TabFM was excluded through an outcome-blind operational amendment because the released implementation was incompatible with the bounded Colab protocol and prohibitively slow.

The completed matrix contains 150 fold tasks and 1,200 method rows. Primary inference is dataset-blocked: folds are averaged within model and the two TFMs are then averaged within dataset. CatBoost is excluded. Thirty dataset--model cells are descriptive only.

\section{Controlled results}\label{sec:controlled}
\subsection{Exact-view reliability mechanisms}
\begin{table}[t]
\caption{Best mean exact expected NLL in each exact-view regime.}
\centering
\begin{tabular}{lll}
\toprule
Regime & Best method & Expected NLL\\
\midrule
Uniform low noise & Geometric pool & 1.478604\\
Direct marginals reliable & Soft CoRe ($\lambda=10$) & 1.486760\\
Conditionals reliable & Arithmetic pool & 1.508830\\
$B\rightarrow A$ reliable & $\Jone$ & 1.475308\\
$A\rightarrow B$ reliable & $\Jtwo$ & 1.475326\\
All four views noisy & Arithmetic pool & 1.538679\\
\bottomrule
\end{tabular}
\end{table}

When direct marginals are reliable, Hard CoRe improves over arithmetic by 0.013270 NLL ($p_{\mathrm{Holm}}=7.11\times10^{-15}$). When conditionals are reliable and direct marginals are noisy, Hard CoRe is worse by 0.150481 with the same adjusted significance order. Exact marginal preservation is therefore reliability-dependent.

\subsection{Heterogeneous task mixture}
Across 100 controlled tasks, Selective CoRe obtains mean exact NLL 1.483684 versus 1.490344 for arithmetic pooling. The paired gain is 0.006660 NLL; Selective wins 68/100 tasks and Wilcoxon $p=1.25\times10^{-11}$. Mean full-family oracle regret is 0.000710. Selection does not collapse to one method: Soft CoRe is selected 38 times, geometric pooling 33, raw orders 19, arithmetic 10, and Hard CoRe zero times.

\subsection{Validation-size sensitivity}
Mean full-family oracle regret decreases from 0.005602 at 100 validation observations to 0.001155 at 800. Mean gain over arithmetic grows from 0.000670 to 0.005118. This supports the qualitative role of validation sample size in flexible policy selection.

\section{Bounded real-model results}\label{sec:real}
\begin{table}[t]
\caption{Dataset-blocked Selective-CoRe effects relative to arithmetic pooling over the two primary TFMs.}
\centering
\begin{tabular}{lrrrr}
\toprule
Metric & Mean difference & 95\% bootstrap CI & Wins & $p$\\
\midrule
Joint NLL & $+0.001831$ & $[+0.000542,+0.003291]$ & 2/10 & 0.02734\\
Joint Brier & $+0.000646$ & $[+0.000288,+0.001059]$ & 1/10 & 0.00586\\
ECE (15 bins) & $+0.000096$ & $[-0.002706,+0.003017]$ & 6/10 & 1.00000\\
Marginal distortion & $-0.003855$ & $[-0.005563,-0.002212]$ & 9/10 & 0.00391\\
\bottomrule
\end{tabular}
\end{table}

Selective CoRe reduces marginal distortion but worsens both proper scores. The joint-NLL disadvantage is positive in eight of ten datasets and remains positive under every leave-one-dataset-out mean. Mean factorization TV is 0.06249 across the 20 primary TFM cells, showing that the negative result is not caused by an absence of incompatibility.

The effect reverses by model. Selective CoRe is worse than arithmetic by 0.006694 NLL on TabICLv2 and better by 0.003031 on TabPFN-3. The paired model-effect difference is 0.009725 ($p=0.00977$; exploratory). Hard CoRe nearly eliminates marginal distortion but is worse than arithmetic by 0.02056 NLL on every primary dataset ($p=0.00195$). The independent joint is worse by 0.13674 NLL. Geometric pooling is descriptively 0.000492 better than arithmetic but not significant ($p=0.7695$).

Pool-only validation selection has mean NLL 1.375591 across the 30 active cells, compared with 1.376046 for the full family, 1.376390 for repair-only, and 1.378435 for raw-only selection. Real validation-fraction sensitivity is non-monotone: 25\%, 50\%, and 100\% of the fixed validation pool yield mean NLL 1.375074, 1.376495, and 1.376046.

\section{Reliability-aware extension protocol}\label{sec:extension}
The audited evidence motivates a second-stage protocol organized around repair opportunity, selection regret, and view reliability. The repository includes a read-only evidence audit script and a dedicated Colab notebook for fresh robustness experiments. Planned runs are: five independent constrained-sampling seeds; context sizes 64/128/256/512/1024; rare-class sensitivity including Customer, Marketing, and Nursery; complexity-aware Safe Selective CoRe with nested families; per-view marginal and conditional reliability metrics; inference-perturbation dispersion; and a third released TFM after operational preflight.

These are explicitly pending experiments. They are not included in the current Results sections until a new evidence package is generated and audited.

\section{Discussion}\label{sec:discussion}
The controlled and real findings identify two regimes rather than a contradiction. Adaptive reconciliation can exploit heterogeneous view reliability when enough validation evidence is available, but a large candidate family can overfit under a 52-example validation budget. Coherence should therefore be viewed as one structural property of a probability system, not an accuracy surrogate.

The model-specific sign reversal is particularly important. It suggests that the usefulness of repair depends on the upstream predictor's reliability structure. Inconsistency magnitude alone cannot identify whether a direct marginal, a conditional direction, or a simple pool deserves trust.

\section{Limitations}\label{sec:limitations}
The current method is pairwise and categorical. The primary real benchmark contains only two genuine TFMs and one deterministic constrained sampling seed. Training context is capped at 256 observations. Multi-seed and context-size robustness are therefore required before a final high-selectivity journal submission.

The audited admission preflight shows extremely rare target classes in Customer (minimum support 3), Marketing (2), and Nursery (2). Rare-class exclusion and support-adaptive sensitivity are needed. The 48-candidate selector is large relative to 52 validation observations and empirically incurs selection cost. ECE is descriptive; conclusions rely primarily on proper scores. Finally, the controlled task mixture is a designed stress test rather than an estimate of the natural prevalence of reliability regimes.

\section{Reproducibility}\label{sec:repro}
The frozen evidence package contains 1,200 method rows across 150 completed fold tasks, controlled replications, ablations, validation sensitivity, fold manifests, data hashes, source revisions, package versions, hardware metadata, plots, and SHA-256 checksums. The archived run used Python 3.13.15, PyTorch 2.11.0 with CUDA 12.8, TabICL 2.1.1, TabPFN 8.4.0, CatBoost 1.2.10, and a Tesla T4. Its frozen source snapshot is commit \texttt{ec6ff45d709096a1665f2475840c466587bc3d26}. Later research-branch changes are not represented as part of that execution snapshot.

A repository-native audit command verifies the expected matrix and recomputes the dataset-blocked primary effect. The 30-cell statistics are descriptive only; confirmatory interpretation uses datasets after averaging the two primary TFMs.

\section{Conclusion}\label{sec:conclusion}
CoRe-TFM shows that probability coherence is not equivalent to probability accuracy. Exact marginal preservation helps only when those marginals deserve trust. Adaptive selection works in controlled heterogeneous reliability regimes with sufficient validation evidence, while the bounded real benchmark shows that a flexible 48-candidate selector can lose to arithmetic pooling even as it reduces structural distortion. The effect also reverses between TabICLv2 and TabPFN-3. The appropriate conclusion is therefore reliability-aware: reconcile only when the available repair opportunity and validation evidence justify the additional complexity.

\bmhead{Acknowledgements}
Funding, full affiliation, postal-address, ORCID, and competing-interest details will be completed before submission. The authors remain responsible for verifying the final evidence and manuscript.

\bibliography{references}
\end{document}
